% GENERATED FILE. Edit canonical lessons, then run npm run generate:books.
% canonical-source-hash: fnv1a64:14e38fef936048cf
% canonical-lessons: GU-C25-savaar, GU-C25-bapor, GU-C25-saanj, GU-C25-divas, GU-C25-mahino, GU-C25-atyaare, GU-C25-time-written, GU-R25-time-written-r1

\chapter{The Time Words Reach the Page}
\label{ch:gu-time-words-written}
\hlchaptermodality{29}

\begin{chapteropening}
\textbf{By the end of this chapter:} \emph{I can read and write the six Gujarati time words I could previously only hear and say.}

Six words that chapter 28 left oral reach the page one at a time, each spelled entirely from signs taught before chapter 8, then a payoff scores reading, speaking, and model-free writing separately and an R1 return writes the other three from sound.
\end{chapteropening}

\section[Practice]{Savār reaches the page}

\label{lesson:GU-C25-savaar}



\begin{quote}
Before anything is uncovered, say the Gujarati word for \textbf{morning}.
\end{quote}

\subsection*{Your turn}
Uncover \textbf{\gu{સવાર}} once. Four pieces, none of them new: \textbf{\gu{સ}} and the \textbf{\gu{ા}} stroke both come from the greeting lesson, and \textbf{\gu{વ}} and \textbf{\gu{ર}} from the courtesy chapter after it. Name the four pieces aloud in order, then read the whole word once.

\subsection*{Writing — guided copy, model in view}
Keep the model in view and copy \textbf{\gu{સવાર}} once. Repair one shape if it needs it; do not start a row.

\begin{tcolorbox}[breakable,colback=teal!4,colframe=teal!35!black,title={Before you move on}]
Which of the four signs in \emph{savār} was taught most recently? (\textbf{\gu{ર}}, in the courtesy chapter — the other three all come out of the greeting.)

Source: \href{https://dictionary.cambridge.org/dictionary/english-gujarati/morning}{Cambridge: morning}.
\end{tcolorbox}
\section[Practice]{Bapor reaches the page}

\label{lesson:GU-C25-bapor}



\begin{quote}
Before anything is uncovered, say the Gujarati word for \textbf{midday}.
\end{quote}

\subsection*{Your turn}
Uncover \textbf{\gu{બપોર}} once. \textbf{\gu{બ}} and \textbf{\gu{પ}} were taught together among the forms that come before your name, the \textbf{\gu{ો}} sign among the doorway nine, and \textbf{\gu{ર}} you wrote yesterday inside \emph{savār}. Name the four pieces aloud in order, then read the whole word once.

\subsection*{Writing — guided copy, model in view}
Keep the model in view and copy \textbf{\gu{બપોર}} once. Repair one shape if it needs it; do not start a row.

\begin{tcolorbox}[breakable,colback=teal!4,colframe=teal!35!black,title={Before you move on}]
Which two signs in \emph{bapor} were taught in the same lesson pair? (\textbf{\gu{બ}} and \textbf{\gu{પ}}, side by side among the forms that come before your name.)

Source: \href{https://dictionary.cambridge.org/dictionary/english-gujarati/afternoon}{Cambridge: midday}.
\end{tcolorbox}
\section[Practice]{Sānj reaches the page}

\label{lesson:GU-C25-saanj}



\begin{quote}
Before anything is uncovered, say the Gujarati word for \textbf{evening}.
\end{quote}

\subsection*{Your turn}
Uncover \textbf{\gu{સાંજ}} once. \textbf{\gu{સ}} and the \textbf{\gu{ા}} stroke opened \emph{savār} two lessons ago; the dot above is the \textbf{\gu{ં}} nasal mark from the doorway nine, and \textbf{\gu{જ}} is the doorway consonant that also ends \emph{āj}. Name the four pieces aloud in order, then read the whole word once.

\subsection*{Writing — guided copy, model in view}
Keep the model in view and copy \textbf{\gu{સાંજ}} once. Repair one shape if it needs it; do not start a row.

\begin{tcolorbox}[breakable,colback=teal!4,colframe=teal!35!black,title={Before you move on}]
Where in \emph{sānj} does the nasal live — at the end, or in the middle? (\textbf{In the middle}, on the \emph{ā}.)

Source: \href{https://dictionary.cambridge.org/dictionary/english-gujarati/evening}{Cambridge: evening}.
\end{tcolorbox}
\section[Practice]{Divas reaches the page}

\label{lesson:GU-C25-divas}



\begin{quote}
Before anything is uncovered, say the Gujarati word for \textbf{day}.
\end{quote}

\subsection*{Your turn}
Uncover \textbf{\gu{દિવસ}} once. Watch the second piece. The \textbf{\gu{િ}} sign is written to the LEFT of \textbf{\gu{દ}} but is read AFTER it, exactly as the forms before your name taught; \textbf{\gu{વ}} and \textbf{\gu{સ}} then follow in order. Name the four pieces aloud in order, then read the whole word once.

\subsection*{Writing — guided copy, model in view}
Keep the model in view and copy \textbf{\gu{દિવસ}} once. Repair one shape if it needs it; do not start a row.

\begin{tcolorbox}[breakable,colback=teal!4,colframe=teal!35!black,title={Before you move on}]
The \textbf{\gu{િ}} sits before \textbf{\gu{દ}} on the page. Which is said first? (\textbf{\gu{દ}} — the sign is written first and read second.)

Source: \href{https://dictionary.cambridge.org/dictionary/english-gujarati/day}{Cambridge: day}.
\end{tcolorbox}
\section[Practice]{Mahino reaches the page}

\label{lesson:GU-C25-mahino}



\begin{quote}
Before anything is uncovered, say the Gujarati word for \textbf{month}.
\end{quote}

\subsection*{Your turn}
Uncover \textbf{\gu{મહિનો}} once. \textbf{\gu{મ}}, \textbf{\gu{હ}} and \textbf{\gu{ન}} all come out of the greeting, the \textbf{\gu{િ}} you used yesterday in \emph{divas}, and the \textbf{\gu{ો}} sign from \emph{bapor}. Name the five pieces aloud in order, then read the whole word once.

\subsection*{Writing — delayed copy, model covered}
Look, cover the model, and write \textbf{\gu{મહિનો}} once. Repair one shape if it needs it; do not start a row.

\begin{tcolorbox}[breakable,colback=teal!4,colframe=teal!35!black,title={Before you move on}]
\emph{Divas} and \emph{mahino} share one vowel sign. Which? (The \textbf{\gu{િ}} — left-written, right-read, in both.)

Source: \href{https://dictionary.cambridge.org/dictionary/english-gujarati/month}{Cambridge: month}.
\end{tcolorbox}
\section[Practice]{Atyāre reaches the page}

\label{lesson:GU-C25-atyaare}



\begin{quote}
Before anything is uncovered, say the Gujarati word for \textbf{now}.
\end{quote}

\subsection*{Your turn}
Uncover \textbf{\gu{અત્યારે}} once. The hard part is in the middle: \textbf{\gu{ત}} carries the \textbf{\gu{્}} stroke that deletes its own vowel and glues it to \textbf{\gu{ય}}, which is exactly the join you first met in \emph{namaste}. Around it stand \textbf{\gu{અ}}, the \textbf{\gu{ા}} stroke, \textbf{\gu{ર}}, and the \textbf{\gu{ે}} sign. Name the six pieces aloud in order, then read the whole word once.

\subsection*{Writing — guided copy, model in view}
Keep the model in view and copy \textbf{\gu{અત્યારે}} once. Repair one shape if it needs it; do not start a row.

\begin{tcolorbox}[breakable,colback=teal!4,colframe=teal!35!black,title={Before you move on}]
Which earlier word taught you the \textbf{\gu{્}} join used here? (\textbf{\gu{નમસ્તે}} — the greeting the book opens with.)

Source: \href{https://dictionary.cambridge.org/dictionary/english-gujarati/now}{Cambridge: now}.
\end{tcolorbox}
\section[Practice]{All eight time words, now on the page}

\label{lesson:GU-C25-time-written}



\begin{quote}
Nothing new arrives here. Say \textbf{morning} and \textbf{now} in Gujarati before any Gujarati is uncovered.
\end{quote}

\subsection*{Your turn}
Two of the four scores are kept here: \textbf{reading} runs page to meaning, \textbf{speaking} runs meaning to sound.

\subsection*{Writing — delayed copy, every model covered}
Write the three from memory, not from a model. Record \textbf{writing} on its own; a 6/6 in reading cannot cover a miss here.

\begin{tcolorbox}[breakable,colback=teal!4,colframe=teal!35!black,title={Before you move on}]
How many of the eight time words can you now write? (\textbf{All eight} — \emph{rāt} and \emph{āj} since the chapter that taught them by ear, and these six today.)

Sources: \href{https://dictionary.cambridge.org/dictionary/english-gujarati/morning}{morning}, \href{https://dictionary.cambridge.org/dictionary/english-gujarati/month}{month}, and \href{https://dictionary.cambridge.org/dictionary/english-gujarati/now}{now}.
\end{tcolorbox}
\section[Practice]{The written time words return at R1}

\label{lesson:GU-R25-time-written-r1}



\begin{quote}
Nothing new arrives here. Say the word for \textbf{now} before anything is read to you.
\end{quote}

\subsection*{Your turn}
\textbf{Listening} runs sound to meaning, exactly as it did before any of these six reached the page.

\subsection*{Writing — from sound, every model covered}
These are the three the payoff did \textbf{not} ask for, so between the two lessons all six have been written without a model.

\begin{tcolorbox}[breakable,colback=teal!4,colframe=teal!35!black,title={Before you move on}]
How far after the payoff did this return come? (\textbf{One lesson} — the first spacing window.)

Sources: \href{https://dictionary.cambridge.org/dictionary/english-gujarati/evening}{evening}, and \href{https://dictionary.cambridge.org/dictionary/english-gujarati/month}{month}.
\end{tcolorbox}
